\documentclass{article}
\usepackage{base}
\usepackage{url}

\title{LaTeX-based Static Website}
\author{Kyler Cain}
\date{2026-07-24}
% tags: latex, webdev

\begin{document}
\maketitle

\begin{abstract}
This site uses a small static publishing system in which every post is written
as a standalone \LaTeX\ document and compiled to both a PDF and HTML version.
\end{abstract}

\section{Objectives}

I wanted to be able to write articles in a familiar format and publish them both
to the web and as a PDF (for printing, sharing, or offline reading) without
maintaining having to think about both formats separately. I also wanted to keep
the site implementation small and simple without a database or application server.

This site is intended to be a personal blog and notes site, not a general-purpose
static-site generator, so I didn't have to think about too many use cases.

\section{Challenges}

My main difficulty is that \LaTeX\ was never intended to produce html output, so
I took a page from arxiv.org\cite{arxiv_accessible_html} and used
LaTeXML\cite{latexml} to handle the conversion.

After getting the basic conversion working, I also wanted to add a tagging system,
a site index, and some basic navigation. This was fairly easy using
Jinja\cite{jinja} templates and a small Python 3\cite{python3} script to read the
\LaTeX\ sources and generate the additional HTML pages.

\section{Implementation}

Each post is a standalone document whose file name is also the URL slug to the post
directory (where I also place the generated PDF). The source files contain a
title, author, date, and abstract that I extract for metadata. The tags are
declared using a comment like \verb|% tags: webdev|.

The builder normalizes each tag into a URL-safe slug and generates a page for each
post and tag. It also generates an index page listing all posts in reverse chronological
order.

The toolchain is packaged in Docker\cite{docker}, and
Makefile\cite{gnu_make} targets handle the build and deployment.
The hosting is done using GitHub Pages, which serves the generated static files directly
from CI/CD build artifacts.

\section{Results}

I have a working static site that is easy to maintain and should be fairly easy to
extend in the future. Adding posts doesn't really require any knowledge of the
publishing system and building/deploying is entirely automated. Hosting via GitHub
Pages is free and simple and the site is fast and reliable thanks to the simplicity
of the static files.

\section{Future Work}

The biggest shortcoming of the current system is that it requires static artifacts for
each post, which means I can't make it build dependencies of posts (such as figures or
output from other scripts).

There is also no paging so you would load the entire index or all of the posts for a tag.
At the scale I am planning to use this site, this isn't a problem (especially since I'm
not incurring any kind of hosting cost), but it may make the site less usable as the
number of posts grows.

\begin{thebibliography}{99}

\bibitem{arxiv_accessible_html}
arXiv.org, ``Accessible HTML,''
\url{https://arxiv.org/help/web_accessibility}

\bibitem{latexml}
Bruce R. Miller, ``LaTeXML: A LaTeX to XML/HTML/MathML Converter,''
\url{https://math.nist.gov/~BMiller/LaTeXML/}

\bibitem{jinja}
Pallets, ``Jinja Documentation,''
\url{https://jinja.palletsprojects.com/en/stable/}

\bibitem{python3}
Python Software Foundation, ``Python 3 Documentation,''
\url{https://docs.python.org/3/}

\bibitem{docker}
Docker, Inc., ``Docker,''
\url{https://docker.com/}

\bibitem{gnu_make}
Free Software Foundation, ``GNU Make Manual,''
\url{https://www.gnu.org/software/make/manual/make.html}

\end{thebibliography}

\end{document}
